\section{Discussion}
\label{sec:discussion}

\subsection{Guidance for monitoring agencies}
The results support five recommendations for remote river monitoring over message-priced, one-way links. They are conditional on the loss models and tolerances studied:
\begin{enumerate}
  \item Send a record within tolerance, not the readings. A four-parameter day needed a median of two 23-byte messages, against at least 47 for the raw readings at 10-minute sampling. Periodic readings and broken-stick points of the same size rarely kept every reading within tolerance and could not be certified.
  \item Protect the whole record and let the node size the code. The record-first code needs no estimate of the loss process and was never worse than any linear design over the same units; repetition is a simpler alternative whose disadvantage is small under long bursts of loss.
  \item Plan for twelve messages per day where loss is bursty. Six messages per day certified most days under independent loss of 0.3 but met the 90\% target at only a few sites under bursty loss; twelve met it at almost every site under both.
  \item Choose the plan from a pilot that includes the active season. A 28-day pilot chose an adequate plan for most site-pilots, but winter pilots underestimated the summer need.
  \item Use certified readings for decisions and report undetermined days as such. Certain decisions were never wrong; near-threshold quantities and diel shape need a tighter tolerance or direct samples.
\end{enumerate}

Neither satellite contact time nor energy was the binding constraint. The module draws 500\,mA at 5\,V for less than one second per message \citep{kineis2023kim1}; therefore, twelve messages per day cost at most 30\,J, and encoding costs approximately a thousandth of that. The constraint was the message quota.

\subsection{Relation to earlier work}
The encoder, the code, and the satellite practice are established \citep{imai1988,reed1960,boehme2009}. The findings qualify their combination for one-way links. The design that is optimal when every message arrives \citep{deligiannakis2007} certified few days once messages could be lost, and a code whose dimension is fixed in advance depended on knowing the loss rate, whereas the record-first code did not. Broken-stick telemetry \citep{boehme2009,photopoulou2015} fixes the number of points rather than the accuracy, which is why it rarely met the tolerance; its pseudo-random transmission was a smaller handicap than its content. Rotating a week of data, as satellite tags do, added time diversity under long bursts at the cost of timeliness, which suggests coding across days as a natural extension. The decision analysis connects the framework to the quality-assurance practice of environmental sensing \citep{campbell2013,horsburgh2015}: a certified record states which readings are trustworthy within a known bound, which flagging alone does not.

\subsection{Limitations}
The results carry the following limitations:
\begin{itemize}
  \item Channels: no measured loss process was available for a satellite IoT service; loss rates and burst lengths were swept, and results are conditional on these synthetic channels. Field measurements could narrow the plan recommendation.
  \item Hardware: no satellite module or microcontroller was used; resource figures for a node are estimates from operation counts. The operator's message format and integrity checks were taken from its public documentation \citep{kineis2024formats} and should be confirmed with the operator.
  \item Integrity: certification assumes that the values used arrived as transmitted. A CRC-16 leaves a residual rate of undetected corruption; the confirmation met it once at a bit-error rate of $10^{-2}$, and the corroboration rule adopted afterwards was evaluated only exploratorily. The operator's device identifier is not cryptographic; a message authentication code would cost one value per message.
  \item Sensor accuracy: the tolerance bounds the error added between the sensor and the server, not the error of the sensor; calibration, fouling, and drift remain the task of quality assurance upstream of the encoder.
  \item Clock: the node clock must stay within half a sampling interval of the true time; a typical crystal requires resynchronisation every few months at 5-minute sampling.
  \item Data: the confirmatory evidence comes from USGS-grade instruments in the United States. Low-cost sensors in Kosovo and Ghana were analysed only exploratorily, on \inwsDays{} and \weijaDays{} days; their records cost more and needed larger plans. Tropical rivers affected by mining, which motivated this study, remain to be evaluated with sufficient data.
  \item Decisions: the criteria and thresholds were illustrative of common regulatory practice; agencies use other criteria, and features finer than the tolerance require a tighter tolerance.
\end{itemize}
